\section{Questão 1}

\subsection{Enunciado}

Um e-commerce possui várias categorias de produtos descritas no conjunto de dados a seguir:
\begin{verbatim}
{
    'Produto': [
        'Eletrônicos',
        'Roupas',
        'Eletrodomésticos',
        'Eletrônicos',
        'Livros',
        'Roupas',
        'Livros'
    ]
}
\end{verbatim}
Aplique a codificação one-hot para representar as categorias de produtos em formato numérico.

\subsection{Solução}

A coluna \texttt{Produto} tem quatro categorias sem ordem: Eletrodomésticos, Eletrônicos, Livros e Roupas. O script monta o \texttt{DataFrame} das sete observações e aplica \texttt{pandas.get\_dummies} com \texttt{dtype=int}. Cada categoria vira uma coluna binária (\texttt{Produto\_<categoria>}): 1 na categoria da linha e 0 nas demais. Cada linha fica com um único 1.

\subsection{Resultado}

A Tabela~\ref{tab:q1-onehot} mostra a matriz. A coluna \texttt{Produto} fica ao lado dos indicadores. Os cabeçalhos omitem o prefixo \texttt{Produto\_}. O arquivo \texttt{scripts/results/q1.csv} guarda as quatro colunas binárias.

\begin{table}[H]
    \centering
    \caption{Codificação one-hot das categorias de produto.}
    \label{tab:q1-onehot}
    \footnotesize
    \setlength{\tabcolsep}{4pt}
    \begin{tabular}{@{}lcccc@{}}
        \toprule
        Produto &
        \shortstack{Eletro-\\domésticos} &
        \shortstack{Eletrô-\\nicos} &
        Livros &
        Roupas \\
        \midrule
        Eletrônicos      & 0 & 1 & 0 & 0 \\
        Roupas           & 0 & 0 & 0 & 1 \\
        Eletrodomésticos & 1 & 0 & 0 & 0 \\
        Eletrônicos      & 0 & 1 & 0 & 0 \\
        Livros           & 0 & 0 & 1 & 0 \\
        Roupas           & 0 & 0 & 0 & 1 \\
        Livros           & 0 & 0 & 1 & 0 \\
        \bottomrule
    \end{tabular}
\end{table}
